% !TeX program = pdflatex
% Illustrative command examples; all values are fictional.
\documentclass{article}
\usepackage[margin=20mm]{geometry}
\usepackage{gradbars,booktabs}
\gradbarsnewcommand[type=bar,args={m O{}},data={#1},options={#2},empty=skip,
 theme=gray,min=0,max=1,width=25mm,height=4mm,precision=3,
 label position=right,label distance=1mm]{\servicebar}
\gradbarsnewcommand[type=bar,args={O{} m},data={#2},options={#1},
 min=0,max=1,width=25mm,height=4mm,precision=3]{\optionfirst}
\gradbarsnewcommand[type=bar,args={m m},data={#1},options={#2},
 min=0,max=1,width=25mm,height=4mm,precision=3]{\requiredbar}
\gradbarsnewcommand[type=dumbbell,args={m m},data={{#1}{#2}},
 min=0,max=1,width=25mm,precision=3]{\pairbar}
\begin{document}
\section*{Custom commands and label placement}
All values are illustrative.
\begin{center}
\renewcommand{\arraystretch}{2.8}
\begin{tabular}{@{}ll@{}}
\toprule
Theme and right label & \servicebar{.960} \\
Solid and left label & \servicebar{.720}[color=teal,label position=left] \\
Gradient and above label & \servicebar{.840}[fill=gradient,left color=blue,right color=cyan!20,label position=above,label reference=value,label yshift=2pt] \\
Below label and offsets & \servicebar{.480}[label position=below,label distance=.4em,label xshift=-2pt] \\
Leading optional options & \optionfirst[color=orange]{.650} \\
Two mandatory arguments & \requiredbar{.530}{color=purple!50} \\
Two observations & \pairbar{.420}{.780} \\
\bottomrule
\end{tabular}
\end{center}
\end{document}
